\clearpage
\section{Appendix: Using the Simulation}

The \href{https://github.com/danielschwalbe/ipres}{simulation} is a project written in the
\href{https://www.python.org}{Python} programming language. In addition to
the \textbf{ipres} Python package, it contains many readmes and example notebooks.
\href{https://jupyter.org}{Jupyter Notebooks} is a technology originating from the Python world
that allows markdown (documentation) and code to be combined in one file and executed
interactively. This makes it suitable for demonstrating such electoral procedures, for example. One can
change parameters and see live how they affect the outcome.

Actually, one would have to start a Jupyter Lab server to execute the notebooks. However, the project
links to \href{https://mybinder.org/v2/gh/danielschwalbe/ipres/HEAD}{binder}, which
copies the project into a virtual machine and already runs the Jupyter server there. One just needs
a bit of patience until the machine has started. Until then, it looks like this:

\begin{figure}[H]
  \centering
  \includegraphics[width=0.85\textwidth]{screenshots/binder_start.png}
  \caption{Binder loading the virtual machine}
\end{figure}

When the machine has booted up, one is greeted by the following screen:

\begin{figure}[H]
  \centering
  \includegraphics[width=0.85\textwidth]{screenshots/binder_loaded.png}
  \caption{JupyterLab after startup}
\end{figure}

To get to the notebooks, one must click on \texttt{notebooks}.

\begin{figure}[H]
  \centering
  \includegraphics[width=0.85\textwidth]{screenshots/binder_loaded_language.png}
   \caption{Notebooks folder}
\end{figure}

\clearpage

And then on \texttt{en} to get to the English notebooks.

\begin{figure}[H]
  \centering
  \clipbox{0pt {\height-7cm} 0pt 0pt} {\includegraphics[width=0.85\textwidth]{screenshots/notebooks_overview.png}}
  \caption{English notebooks}
\end{figure}

Here one can then pick one out, e.g., \texttt{federal\_election.ipynb}.

By clicking on the start arrow, the currently selected cell is executed. Note that the
cells generally must be executed from top to bottom, because the lower cells
usually (not always) require the results of the upper cells.

\begin{figure}[H]
  \centering
   \clipbox{0pt {\height-7cm} 0pt 0pt}{\includegraphics[width=0.85\textwidth]{screenshots/example_run_cell.png}}
  \caption{Execution of a notebook cell}
\end{figure}

\clearpage
Using the two arrows, one can also execute the entire notebook at once:

\begin{figure}[H]
  \centering
  \clipbox{0pt {\height-7cm} 0pt 0pt}{\includegraphics[width=0.85\textwidth]{screenshots/example_run_all.png}}
  \caption{Execute all cells at once}
\end{figure}

Then one will be asked whether to restart the kernel.

\begin{figure}[H]
  \centering
  \includegraphics[width=0.85\textwidth]{screenshots/restart_kernel.png}
  \caption{Restart kernel}
\end{figure}

As long as one has not changed anything in the notebook, one can always say yes here. Only if through continued experimentation
one has brought the variables into a state that cannot be reconstructed from the notebook code, must one know that
the result will be gone after the kernel restart.

If one wants to change something in the code, one can simply edit the lines and execute them again.
One could, for example, enter different seeds in the upper cell and would see how other parties
win.

If one wants to make more extensive changes, one does not necessarily need to
know Python nowadays. AIs can do that too and are happy to help.
